% =============================================================================
%  Préambule du guide technique InsaCloud
% =============================================================================
\usepackage{fontspec}
% Noms de fichiers plutôt que noms système : la police vient du paquet
% tectonic, le document se recompile donc à l'identique sur n'importe quel poste.
\setmainfont{lmroman10-regular.otf}[
  BoldFont=lmroman10-bold.otf, ItalicFont=lmroman10-italic.otf,
  BoldItalicFont=lmroman10-bolditalic.otf]
\setsansfont{lmsans10-regular.otf}[
  BoldFont=lmsans10-bold.otf, ItalicFont=lmsans10-oblique.otf]
\setmonofont{lmmono10-regular.otf}[
  BoldFont=lmmonolt10-bold.otf, ItalicFont=lmmono10-italic.otf, Scale=0.92]
\usepackage[french]{babel}
\usepackage{graphicx,xcolor,booktabs,array,tabularx,longtable,multirow}
\usepackage{tikz}
\usetikzlibrary{positioning,arrows.meta,shapes.geometric,calc,fit,backgrounds,
                decorations.pathreplacing,patterns,shadows.blur}
\usepackage{fancyhdr,enumitem,microtype}
\usepackage[colorlinks=true,linkcolor=insared,urlcolor=insared,citecolor=insared]{hyperref}
\usepackage{caption}
\usepackage{listings}
\usepackage{tcolorbox}
\tcbuselibrary{skins,breakable}
\usepackage{titlesec}
\usepackage{emptypage}

% --- Couleurs ----------------------------------------------------------------
\definecolor{insared}{RGB}{200,16,46}
\definecolor{grisfond}{RGB}{247,248,250}
\definecolor{grisligne}{RGB}{216,220,226}
\definecolor{gristexte}{RGB}{91,97,107}
\definecolor{vertok}{RGB}{26,127,75}
\definecolor{vertfond}{RGB}{238,247,241}
\definecolor{orangepartiel}{RGB}{178,106,5}
\definecolor{orangefond}{RGB}{254,246,231}
\definecolor{rougeko}{RGB}{198,40,40}
\definecolor{bleufond}{RGB}{238,243,250}
\definecolor{bleutrait}{RGB}{43,91,158}

% --- Mise en page ------------------------------------------------------------
\captionsetup{font=small,labelfont={bf,color=insared}}
\setlist[itemize]{leftmargin=1.3em,itemsep=2pt,topsep=3pt,parsep=0pt}
\setlist[enumerate]{leftmargin=1.5em,itemsep=2pt,topsep=3pt,parsep=0pt}
\setlist[description]{leftmargin=1.3em,itemsep=3pt,topsep=3pt,parsep=0pt}

% LaTeX refuse par defaut qu'une figure occupe plus de 70 % d'une page et
% exige 20 % de texte : avec beaucoup de schemas, cela repousse les flottants
% et laisse des pages a moitie vides. On relache ces seuils.
\renewcommand{\topfraction}{0.92}
\renewcommand{\bottomfraction}{0.85}
\renewcommand{\textfraction}{0.06}
\renewcommand{\floatpagefraction}{0.75}
\setcounter{topnumber}{3}
\setcounter{bottomnumber}{2}
\setcounter{totalnumber}{4}

\pagestyle{fancy}
\fancyhf{}
\fancyhead[L]{\small\textcolor{gristexte}{\leftmark}}
\fancyhead[R]{\small\textcolor{gristexte}{InsaCloud}}
\fancyfoot[C]{\small\thepage}
\renewcommand{\headrulewidth}{0.4pt}
\renewcommand{\headrule}{\hbox to\headwidth{\color{grisligne}\leaders\hrule height \headrulewidth\hfill}}

\titleformat{\chapter}[display]
  {\normalfont\huge\bfseries}
  {\textcolor{insared}{\normalsize\MakeUppercase{\chaptertitlename} \thechapter}}
  {6pt}{\Huge}
\titlespacing*{\chapter}{0pt}{-20pt}{24pt}

% --- Extraits de code --------------------------------------------------------
\lstdefinestyle{insacloud}{
  basicstyle=\ttfamily\footnotesize,
  backgroundcolor=\color{grisfond},
  frame=single, rulecolor=\color{grisligne},
  breaklines=true, breakatwhitespace=false,
  postbreak=\mbox{\textcolor{gristexte}{$\hookrightarrow$}\space},
  showstringspaces=false,
  keywordstyle=\color{insared}\bfseries,
  commentstyle=\color{gristexte}\itshape,
  stringstyle=\color{vertok},
  numberstyle=\tiny\color{gristexte},
  xleftmargin=6pt, framexleftmargin=6pt, aboveskip=8pt, belowskip=8pt,
  extendedchars=true, inputencoding=utf8,
  literate=%
    {é}{{\'e}}1 {è}{{\`e}}1 {ê}{{\^e}}1 {ë}{{\"e}}1
    {à}{{\`a}}1 {â}{{\^a}}1 {ä}{{\"a}}1
    {î}{{\^i}}1 {ï}{{\"i}}1
    {ô}{{\^o}}1 {ö}{{\"o}}1
    {û}{{\^u}}1 {ù}{{\`u}}1 {ü}{{\"u}}1
    {ç}{{\c c}}1 {É}{{\'E}}1 {È}{{\`E}}1 {À}{{\`A}}1 {Ç}{{\c C}}1
    {œ}{{\oe}}1 {Œ}{{\OE}}1 {æ}{{\ae}}1
    {—}{{---}}1 {–}{{--}}1 {’}{{'}}1 {‘}{{'}}1
    {«}{{\guillemotleft}}1 {»}{{\guillemotright}}1
    {·}{{\textperiodcentered}}1 {→}{{$\rightarrow$}}1 {←}{{$\leftarrow$}}1
    {✔}{{\checkmark}}1 {°}{{\textdegree}}1 {€}{{\texteuro}}1
}
\lstset{style=insacloud}

\lstdefinelanguage{yamlplus}{
  keywords={true,false,null,name,hosts,become,roles,tasks,vars,when,loop,with_items,
            state,mode,owner,group,path,src,dest,notify,handlers,services,image,
            environment,volumes,networks,ports,depends_on,build,command,restart},
  keywordstyle=\color{insared}\bfseries,
  sensitive=false,
  comment=[l]{\#},
  morestring=[b]",
  morestring=[b]',
}

\lstdefinelanguage{dockerfile}{
  keywords={FROM,RUN,CMD,LABEL,EXPOSE,ENV,ADD,COPY,ENTRYPOINT,VOLUME,USER,
            WORKDIR,ARG,ONBUILD,STOPSIGNAL,HEALTHCHECK,SHELL,AS},
  keywordstyle=\color{insared}\bfseries,
  sensitive=true,
  comment=[l]{\#},
  morestring=[b]",
}

% --- Encadrés pédagogiques ---------------------------------------------------
\newtcolorbox{notion}[1]{
  breakable, enhanced, colback=bleufond, colframe=bleutrait,
  boxrule=0pt, leftrule=3pt, arc=2pt, left=8pt, right=8pt, top=6pt, bottom=6pt,
  fonttitle=\bfseries\small, title={Notion — #1},
  coltitle=bleutrait, titlerule=0pt, colbacktitle=bleufond,
  attach title to upper=\par\smallskip
}

\newtcolorbox{piege}[1]{
  breakable, enhanced, colback=orangefond, colframe=orangepartiel,
  boxrule=0pt, leftrule=3pt, arc=2pt, left=8pt, right=8pt, top=6pt, bottom=6pt,
  fonttitle=\bfseries\small, title={Piège — #1},
  coltitle=orangepartiel, titlerule=0pt, colbacktitle=orangefond,
  attach title to upper=\par\smallskip
}

\newtcolorbox{retenir}{
  breakable, enhanced, colback=vertfond, colframe=vertok,
  boxrule=0pt, leftrule=3pt, arc=2pt, left=8pt, right=8pt, top=6pt, bottom=6pt,
  fonttitle=\bfseries\small, title={À retenir},
  coltitle=vertok, titlerule=0pt, colbacktitle=vertfond,
  attach title to upper=\par\smallskip
}

% --- Raccourcis --------------------------------------------------------------
\newcommand{\ok}{\textcolor{vertok}{\textbf{Conforme}}}
\newcommand{\partiel}{\textcolor{orangepartiel}{\textbf{Partiel}}}
\newcommand{\plus}{\textcolor{insared}{\textbf{Au-delà}}}
\newcommand{\fichier}[1]{\textcolor{insared}{\texttt{#1}}}
\newcommand{\cmd}[1]{\texttt{#1}}

% Les paragraphes techniques enchaînent de longues suites en \texttt{} que TeX
% ne peut pas couper : sans marge de manœuvre, elles débordent dans la marge.
\emergencystretch=3em
\hbadness=6000
\tolerance=2000

% =============================================================================
%  Propre au document « technologies » : une fiche par outil
% =============================================================================
\usepackage[export]{adjustbox}

% Un bandeau d'ouverture pour chaque outil, avec son numéro de brique.
\newcommand{\outil}[2]{%
  \vspace{10pt}%
  \noindent\begin{tcolorbox}[enhanced, colback=grisfond, colframe=insared,
    boxrule=0pt, leftrule=4pt, arc=2pt, left=10pt, right=10pt,
    top=7pt, bottom=7pt, breakable]
  {\Large\bfseries #1}\hfill{\small\color{gristexte}#2}
  \end{tcolorbox}\vspace{2pt}}

% Les quatre questions auxquelles chaque fiche répond.
\newcommand{\quoi}{\par\vspace{5pt}\noindent{\color{insared}\textbf{C'est quoi}}\par\nopagebreak}
\newcommand{\pourquoi}{\par\vspace{7pt}\noindent{\color{insared}\textbf{Pourquoi celui-là}}\par\nopagebreak}
\newcommand{\role}{\par\vspace{7pt}\noindent{\color{insared}\textbf{Son rôle ici}}\par\nopagebreak}
\newcommand{\apport}{\par\vspace{7pt}\noindent{\color{insared}\textbf{Ce que ça a changé}}\par\nopagebreak}

% Capture recadrée : on ne montre que la zone dont parle le texte.
\newcommand{\gros}[3][0.8]{%
  \begin{figure}[htbp]\centering
  \includegraphics[width=#1\textwidth,frame]{figures/#2}
  \caption{#3}\end{figure}}

\renewcommand{\topfraction}{0.92}
\renewcommand{\bottomfraction}{0.85}
\renewcommand{\textfraction}{0.06}
\renewcommand{\floatpagefraction}{0.75}
